% 27

\eocesol{(a)~Geometrik, 0.0804.
(b)~Binomial, 0.0322.
(c)~Binomial negatif, 0.0193.}

% 29

\eocesol{(a)~Binomial negatif dengan $n=4$ dan $p=0.55$, dengan sukses didefinisikan sebagai siswa perempuan. Kasus binomial negatif tepat digunakan karena hasil percobaan terakhir ditetapkan, sedangkan urutan 3 percobaan pertama tidak diketahui.
(b)~0.1838.
(c)~${3 \choose 1} = 3$.
(d)~Dalam model binomial tidak ada pembatasan pada hasil percobaan terakhir. Dalam model binomial negatif, hasil percobaan terakhir ditetapkan. Oleh karena itu, kita memperhatikan banyaknya cara mengurutkan $k - 1$ sukses lainnya dalam $n - 1$ percobaan pertama.}
